% !PRINTED 689
\weberpage{689}

Die Bezeichnung (5) eines Differentials unterscheidet sich von
der ähnlichen Bezeichnung der \gesperrt{Funktionen} in $\Omega$ dadurch, daß
der Grad des Zählers um $(2p - 2)$ höher ist als der des Nenners.

Die in § 194 definierte Differentiale $d\alpha$ der Funktionen in
% [sic] printed "definierte", the plural "Die ... Differentiale ... sind" requires "definierten"
$\Omega$ sind spezielle Fälle der allgemeinen Differentiale (5). Denn
nicht alle diese $dJ$ sind Differentiale von Funktionen in $\Omega$.
Wir unterscheiden \gesperrt{eigentliche Differentiale} $d\alpha$, d. h. solche,
die Differentiale von Funktionen in $\Omega$ sind, und \gesperrt{uneigentliche}
oder \gesperrt{Abelsche Differentiale} $dJ$, die das nicht sind. Für die
letzteren hat $J$ selbst keine Bedeutung und erhält eine solche
erst in der Integralrechnung, die der rein arithmetischen Methode
nicht zugänglich ist.

Ist $z$ eine beliebige Variable in $\Omega$ und $dJ$ ein Differential,
so ist $dJ:dz$ eine Funktion in $\alpha$,
% [sic] printed $\alpha$, a misprint for $\Omega$ (dJ:dz is a function in the field; cf. "gleichfalls immer eine Funktion in $\Omega$" two sentences later)
und $dJ:dz$ ist ein \gesperrt{Diffe\-rentialquotient}.
Wir unterscheiden auch hier \gesperrt{eigentliche und uneigentliche}
Differentialquotienten. Der Quotient zweier Differentiale $dJ:dJ'$
ist gleichfalls immer eine Funktion in $\Omega$ und kann ebenfalls
ein Differentialquotient genannt werden.

Damit ein Polygonquotient $\mathfrak{A}:\mathfrak{B}$ das Symbol für ein Differential
sein kann, ist notwendig, daß die Differenz $(a - b)$ der Ordnungen von
$\mathfrak{A}$ und $\mathfrak{B}$ gleich $(2p - 2)$ sei. Aber diese Bedingung ist nicht
hinreichend; es kommt noch hinzu, daß eine Variable $z$ in $\Omega$
existieren muß, für die
\weq{6}{\mathfrak{U}^2\mathfrak{A} \sim \mathfrak{Z}\mathfrak{B}}
ist. Ist diese Forderung befriedigt, so bleibt sie erhalten, wenn
$z$ durch irgend eine andere Variable in $\Omega$ und $\mathfrak{A}$ und $\mathfrak{B}$ durch
äquivalente Polygone ersetzt werden. Sind also $\mathfrak{A}'$, $\mathfrak{A}''$, \dots{} Poly\-gone
der durch $\mathfrak{A}$ bestimmten Klasse $A$, so sind $\mathfrak{A}':\mathfrak{B}$, $\mathfrak{A}'':\mathfrak{B}$, \dots{}
gleichfalls Differentiale.

Ist
\[ A = (\mathfrak{A}_1,\ \mathfrak{A}_2,\ \mathfrak{A}_3,\ \dots) \]
und $\mathfrak{A}_1$, $\mathfrak{A}_2$, $\mathfrak{A}_3$, \dots{} eine Basis von $A$, so sind die entsprechenden
Differentialquotienten
\[ \frac{dJ_1}{dz},\ \frac{dJ_2}{dz},\ \frac{dJ_3}{dz},\ \dots \]
die Basis einer Funktionenschar von endlicher Dimension, und wir
können
\[ (dJ_1,\ dJ_2,\ dJ_3,\ \dots) \]
die Basis einer Schar von Differentialen von derselben Dimension
nennen.
